\documentclass[12pt]{article}

\input{../../_assets/preambulo.tex}


\begin{document}

    % 1. Foto de fondo
    % 2. Título
    % 3. Encabezado Izquierdo
    % 4. Color de fondo
    % 5. Coord x del titulo
    % 6. Coord y del titulo
    % 7. Fecha

    
    \input{../../_assets/portada}
    \portadaExamen{ffccA4.jpg}{Ecuaciones\\Diferenciales II\\Examen VII}{Ecuaciones Diferenciales II. Examen VII}{MidnightBlue}{-8}{28}{2026}{}

    \begin{description}
        \item[Asignatura] Ecuaciones Diferenciales II.
        \item[Curso Académico] 2014/2015.
        \item[Grado] Doble Grado en Ingeniería Informática y Matemáticas.
        \item[Grupo] Único.
        \item[Profesor] Rafael Ortega Ríos.
        \item[Descripción] Segundo examen.
        \item[Fecha] 9 de junio de 2015.
        % \item[Duración] ---.
    \end{description}
    \newpage


    % ------------------------------------
    
    \begin{ejercicio}
        Encuentra la solución general de la ecuación diferencial $\dot{x} = tx^3$. En particular se especificará el dominio de definición $\cc{D} \subset \bb{R}^2$.
    \end{ejercicio}

    \begin{ejercicio}
        Calcula un número real $\mu$ para el que se cumpla $\lim\limits_{t \to +\infty} \frac{\|e^{At}\|}{e^{\mu t}} = 0$ si $A$ es la matriz
        \[
        \begin{pmatrix}
            3 & 0 & 0 \\
            0 & 5 & 4 \\
            0 & -4 & 5
        \end{pmatrix}
        \]
    \end{ejercicio}

    \begin{ejercicio}
        La solución del problema de valores iniciales
        \[
        \ddot{x} + x + x^3 = 0, \quad x(0) = 2\alpha, \quad \dot{x}(0) = \alpha
        \]
        se denota por $x(t; \alpha)$. Calcula $x(t; 0)$ y $\del{x}{\alpha}(t; 0)$ para cada $t \in \bb{R}$. Explica los teoremas que se emplean para justificar el cálculo anterior.
    \end{ejercicio}

    \begin{ejercicio}
        Se considera una sucesión de funciones continuas $f_n : \left[0, 1\right] \to \bb{R}$ que converge uniformemente a $f : \left[0, 1\right] \to \bb{R}$ y una sucesión de números $t_n \in \left[0, 1\right]$ que converge a $t_*$. Demuestra que $f_n(t_n) \to f(t_*)$.
    \end{ejercicio}

    \begin{ejercicio}
        Se considera el problema de valores iniciales
        \[
        \dot{x} = \frac{f(t,x)}{t^2 + x^2}, \quad x(0) = 1
        \]
        donde $f : \bb{R}^2 \to \bb{R}$ es continua y cumple la condición $|f(t,x)| \leq 2|x| + 7$ para cada $(t,x) \in \bb{R}^2$. Demuestra que existe una solución con intervalo maximal $\left]-\infty, +\infty\right[$.
    \end{ejercicio}

    \newpage
    \setcounter{ejercicio}{0} % Reiniciar contador de ejercicios
    % \noindent
    % \textbf{Solución.}

\end{document}